\section{引言：从3D重建到生成式3D}

Philipp Henzler 自2018年攻读博士期间便开始从事3D重建与生成方向的研究，主要工作包括自监督3D重建（self-supervised 3D reconstruction）和外观建模（appearance modeling）。2022年加入 Google 后，他致力于生成式3D重建和3D重光照（relighting），其团队成功将研究成果应用于 Google Shopping 产品中——先从鞋类3D重建开始，后借助视频模型扩展到任意物体品类。

本次演讲的核心主题是：\textbf{3D生成方法如何与通用视频模型相结合}，以及在视频模型日益强大的今天，显式3D建模为何依然不可或缺。

\begin{importantbox}{演讲核心问题}
在2D视频模型已经能够生成高度逼真视频的今天，我们为什么还需要显式建模3D？3D方法与2D生成模型各自的优势和劣势是什么？如何将二者结合以获得最佳效果？
\end{importantbox}

